% ============================================================================
% chapters/ch3_methodology.tex
% ============================================================================
\chapter{Methodology}
\label{ch:methodology}

% ============================================================================
\section{System Overview}
\label{sec:system_overview}
% ============================================================================

The proposed system is a stage-based production pipeline that transforms raw
wearable sensor recordings into activity predictions and continuously monitors,
evaluates, and adapts the underlying model.
The pipeline comprises 14 stages, logically partitioned into three groups that
can be executed independently or in combination.

\begin{table}[ht]
  \centering
  \caption{Pipeline stage groups and their purposes.}
  \label{tab:stage_groups}
  \begin{tabularx}{\linewidth}{llX}
    \toprule
    \textbf{Group} & \textbf{Stages} & \textbf{Purpose} \\
    \midrule
    Inference Cycle     & 1--7   & Ingest data, predict activities, monitor quality \\
    Retraining Cycle    & 8--10  & Retrain, register, and update baselines \\
    Advanced Analytics  & 11--14 & Calibrate, detect drift, pseudo-label, adapt \\
    \bottomrule
  \end{tabularx}
\end{table}

This decomposition follows the \textit{composable pipeline} design principle:
each stage is an independently executable unit with a typed configuration input
and a typed artefact output.
The orchestrator composes stages at runtime according to command-line arguments.
Default execution runs stages 1--7 only; retraining and advanced analytics are
enabled via explicit flags, ensuring expensive computations occur only when needed.

\subsection{Complete Stage Enumeration}

\begin{table}[ht]
  \centering
  \caption{All 14 pipeline stages with inputs and outputs.}
  \label{tab:all_stages}
  \begin{tabularx}{\linewidth}{clXX}
    \toprule
    \textbf{\#} & \textbf{Stage Name} & \textbf{Input} & \textbf{Output} \\
    \midrule
    1  & Data Ingestion          & Raw Excel / CSV files          & Fused resampled CSV (50 Hz) \\
    2  & Data Validation         & Fused CSV                      & Validation report (pass/fail) \\
    3  & Data Transformation     & Validated CSV                  & Windowed NumPy arrays (.npy) \\
    4  & Model Inference         & .npy arrays + Keras model      & Prediction CSV + probabilities \\
    5  & Model Evaluation        & Predictions + probabilities    & Confidence distribution, ECE \\
    6  & Post-Inference Monitor  & Predictions (unlabelled)       & Three-layer monitoring report \\
    7  & Trigger Evaluation      & Monitoring report              & Retrain decision + reason \\
    8  & Model Retraining        & Training data + current model  & Retrained model checkpoint \\
    9  & Model Registration      & Retrained model + metrics      & Versioned model in registry \\
    10 & Baseline Update         & New model + reference data     & Updated drift baselines \\
    11 & Calibration \& UQ       & Validation logits              & Temperature $T$, ECE, uncertainty \\
    12 & Wasserstein Drift       & Production features + baseline & Per-channel drift scores \\
    13 & Curriculum Pseudo-Label & Model + unlabelled data        & Fine-tuned model \\
    14 & Sensor Placement        & 6-axis features                & Hand-side label, mirrored features \\
    \bottomrule
  \end{tabularx}
\end{table}

% ============================================================================
\section{Dataset and Preprocessing}
\label{sec:dataset}
% ============================================================================

\subsection{Data Collection}

The dataset comprises 26 wrist sensor recording sessions collected from
Garmin smartwatches in Decoded CSV export format.
Each session produces three CSV files: an accelerometer file
(\texttt{\{ts\}\_accelerometer.csv}), a gyroscope file
(\texttt{\{ts\}\_gyroscope.csv}), and a metadata file
(\texttt{\{ts\}\_record.csv}).
Sessions were originally recorded at approximately 50 Hz with a target
resampling rate of exactly 50 Hz enforced during ingestion.

\subsection{Sensor Fusion and Resampling}

The ingestion stage merges the accelerometer and gyroscope streams on their
nearest timestamps (nearest-neighbour temporal alignment), producing a unified
6-channel array with columns
$[\,\text{acc}_x,\,\text{acc}_y,\,\text{acc}_z,\,
   \text{gyro}_x,\,\text{gyro}_y,\,\text{gyro}_z\,]$.
Re-sampling to exactly 50 Hz is performed using linear interpolation.

\paragraph{Unit normalisation.}
A critical preprocessing step, motivated by a discovered production incident
(see \Cref{sec:unit_mismatch}), is automatic unit detection.
The pipeline computes the median magnitude
$\|\mathbf{a}\| = \sqrt{a_x^2 + a_y^2 + a_z^2}$ over the first 100 windows.
If the median exceeds 100\,m/s\textsuperscript{2}, the data are assumed to be
in milli-g (1\,g~$\approx$~9807\,milli-g) and are converted:
\begin{equation}
  \mathbf{a}^{(\text{SI})} = \mathbf{a}^{(\text{mg})} \times
    \frac{9.80665}{1000}
\end{equation}
This guard ensures forward compatibility with future Garmin firmware versions
that may change the export unit.

\subsection{Windowing}

Following \citet{bao2004activity} and \citet{chen2021deep}, raw time-series
data are segmented into fixed-length windows of $W = 200$ samples (4 seconds at
50 Hz) with 50\,\% overlap, yielding a stride of 100 samples per window.
Each window is normalised per channel using a StandardScaler fitted on the
training set.
The resulting tensors have shape
$(\text{windows},\,W,\,6) = (\text{windows},\,200,\,6)$.

% ============================================================================
\section{Model Architecture}
\label{sec:model_architecture}
% ============================================================================

The classification model is a \cnnbilstm{} with the following topology:

\begin{table}[ht]
  \centering
  \caption{\cnnbilstm{} architecture. Total parameters: 499,131.}
  \label{tab:model_architecture}
  \begin{tabular}{clcc}
    \toprule
    \textbf{Layer} & \textbf{Type} & \textbf{Units/Filters} & \textbf{Output shape} \\
    \midrule
    1  & Input            & ---       & (200, 6) \\
    2  & Conv1D + BN      & 64        & (200, 64) \\
    3  & Conv1D + BN      & 64        & (200, 64) \\
    4  & MaxPool + Drop (0.1)  & ---  & (100, 64) \\
    5  & Conv1D + BN      & 128       & (100, 128) \\
    6  & Conv1D + BN      & 128       & (100, 128) \\
    7  & MaxPool + Drop (0.2)  & ---  & (50, 128) \\
    8  & BiLSTM           & 128       & (50, 256) \\
    9  & BiLSTM           & 64        & (128) \\
    10 & Dense + BN       & 256       & (256) \\
    11 & Dropout (0.3)    & ---       & (256) \\
    12 & Dense (softmax)  & 11        & (11) \\
    \bottomrule
  \end{tabular}
\end{table}

The model was pre-trained with the Adam optimiser \citep{kingma2015adam}
($\text{lr} = 10^{-3}$, $\beta_1 = 0.9$, $\beta_2 = 0.999$),
batch size 32, and 5-fold cross-validation \citep{kohavi1995study}.
The best fold achieved validation accuracy of 89.9\,\%.

% ============================================================================
\section{Three-Layer Monitoring Framework}
\label{sec:monitoring}
% ============================================================================

Post-inference monitoring operates entirely on unlabelled production
predictions.
Three complementary layers are computed after each inference batch.

\subsection{Layer 1: Confidence Monitoring}

Layer 1 computes three statistics from the per-window softmax probability vector:
\begin{itemize}
  \item \textbf{Mean confidence} $\bar{c}$: average maximum softmax probability
        across all windows in the batch.
  \item \textbf{Uncertain window fraction} $u$: fraction of windows with
        $\max_k p_k < \theta_{\text{warn}} = 0.60$.
  \item \textbf{Very uncertain fraction}: fraction with
        $\max_k p_k < \theta_{\text{critical}} = 0.40$.
\end{itemize}
A warning is raised when $u > 30\,\%$; a critical alert is raised when
$u > 50\,\%$.

\subsection{Layer 2: Temporal Consistency}

Layer 2 measures the coherence of the predicted activity sequence over time:
\begin{itemize}
  \item \textbf{Transition rate} $r_t$: fraction of consecutive window pairs
        where the predicted class changes.
        Alert when $r_t > 50\,\%$ (erratic predictions, likely misaligned sensor).
  \item \textbf{Mean dwell time} $\bar{d}$: average number of consecutive
        windows sharing the same predicted class.
        Alert when $\bar{d} < 3$ windows (< 0.6 s at 50\,\% overlap).
\end{itemize}

\subsection{Layer 3: Distribution Drift}

Layer 3 computes per-channel feature \textit{z}-scores relative to the
reference baseline stored in \texttt{models/normalized\_baseline.json}:
\begin{equation}
  z_c = \frac{\mu_c^{\text{prod}} - \mu_c^{\text{ref}}}
             {\sigma_c^{\text{ref}}}
\end{equation}
where $c$ indexes the six sensor channels.
A \textit{drift warning} is issued when $|z_c| > 2.0$ for any channel.
The Population Stability Index (from \Cref{eq:psi}) is computed over the full
24-bin histogram; the multi-channel aggregated PSI uses warn/critical thresholds
of 0.75 / 1.50 (empirically calibrated in \Cref{sec:psi_calibration}).

% ============================================================================
\section{Trigger Policy Engine}
\label{sec:trigger}
% ============================================================================

The trigger engine translates monitoring signals into a binary retraining
decision.
It implements a \textit{two-of-three voting rule}: a retraining trigger is
issued only when at least two of the three monitoring layers report a warning
or critical state.
This prevents a single noisy signal from causing spurious retraining.

The trigger logic also enforces:
\begin{itemize}
  \item \textbf{Cooldown period}: after a trigger fires, no further trigger
        can be issued for 24 hours, preventing oscillatory retraining.
  \item \textbf{Tiered severity}: the trigger decision distinguishes
        \textsc{info}, \textsc{warning}, and \textsc{critical} levels.
        Only \textsc{warning} and \textsc{critical} initiate retraining.
  \item \textbf{Forced retrain flag}: the \texttt{forced\_retrain} parameter
        bypasses the voting rule for scheduled maintenance windows.
\end{itemize}

% ============================================================================
\section{Adaptation Strategies}
\label{sec:adaptation}
% ============================================================================

\subsection{AdaBN + TENT}
\label{sec:tent_adabn}

The first adaptation strategy chains AdaBN \citep{li2018adabn} and TENT
\citep{wang2021tent}:

\begin{enumerate}
  \item \textbf{AdaBN phase.} Run $n_{\text{batches}} = 10$ forward passes over
        the target domain data in training mode.
        After each pass, the Batch Normalisation running statistics are updated
        to reflect the target domain.
  \item \textbf{Snapshot.} Before the TENT gradient loop begins, record a
        snapshot of all BN layers' \texttt{running\_mean} and
        \texttt{running\_var} tensors.
  \item \textbf{TENT phase.} For each of $n_{\text{steps}} = 10$ gradient steps
        at $\text{lr} = 10^{-4}$, minimise the entropy of the output softmax
        distribution via gradient descent on $\gamma$ and $\beta$ only.
        After \textit{each} gradient step, restore the BN running statistics
        from the snapshot.
        This prevents TENT's training-mode forward passes from corrupting the
        statistics set by AdaBN.
  \item \textbf{Safety gates.}
        (a)~\textit{OOD gate}: if the initial mean normalised entropy of the
        target batch exceeds $\theta_{\text{ood}} = 0.85$, skip TENT and return
        the AdaBN-adapted model unchanged.
        (b)~\textit{Confidence rollback}: if mean confidence drops by more than
        1 percentage point after TENT relative to the AdaBN checkpoint,
        restore all $\gamma, \beta$ to their pre-TENT values.
\end{enumerate}

\subsection{Curriculum Pseudo-Labelling}
\label{sec:pseudo_labelling}

The second adaptation strategy fine-tunes the model on target-domain data
using pseudo-labels generated by the model itself:

\begin{enumerate}
  \item \textbf{Confidence threshold initialisation.} Set initial threshold
        $\tau_0 = 0.95$.
  \item For each curriculum iteration $i$:
    \begin{enumerate}[label=\alph*.]
      \item Generate predictions on the unlabelled target set.
      \item Retain windows with $\max_k p_k > \tau_i$ as pseudo-labelled data.
      \item Apply \textit{self-consistency filter}: run a second forward pass;
            retain only windows where the top-1 prediction is identical in both
            passes.
      \item Fine-tune the model on the retained pseudo-labelled windows with EWC:
            \[
              \mathcal{L} = \mathcal{L}_{\text{CE}} +
                \frac{\lambda}{2} \sum_j F_j (\theta_j - \theta_j^*)^2
            \]
            where $F_j$ is the Fisher information for parameter $j$,
            $\theta_j^*$ is the pre-training value, and $\lambda = 1000$
            \citep{kirkpatrick2017overcoming}.
      \item Lower the threshold: $\tau_{i+1} = \tau_i - 0.05$.
    \end{enumerate}
  \item \textbf{Rollback guard.} Evaluate on a held-out source-domain
        validation set before and after each iteration.
        If validation accuracy drops by more than 10 percentage points,
        roll back to the previous checkpoint.
\end{enumerate}

The source-domain data used for rollback evaluation is pre-processed with
the production StandardScaler before being mixed with pseudo-labelled data.
This alignment step was introduced after a problem in which source data was
passed at raw scale, inflating in-loop validation accuracy while degrading
production performance (see \Cref{sec:engineering_challenges}).

% ============================================================================
\section{Calibration and Uncertainty Quantification}
\label{sec:calibration}
% ============================================================================

\subsection{Temperature Scaling}

Temperature scaling optimises a scalar temperature $T$ by minimising ECE on
a held-out validation set.
The initial value is $T_0 = 1.5$ \citep{guo2017calibration}.
The temperature is optimised over 50 Adam steps with learning rate $10^{-2}$.
A warning is issued when post-calibration ECE exceeds 0.10
\citep{naeini2015obtaining}.

\subsection{Monte Carlo Dropout}

At inference time, $T = 30$ stochastic forward passes are run with dropout
active \citep{gal2016dropout}.
The mean prediction is used as the point estimate; the variance across passes
is the epistemic uncertainty estimate.
Windows with mean epistemic uncertainty above the 90th percentile of the
training distribution are flagged as candidates for active learning export.
